\documentclass[conference]{IEEEtran}

% Essential packages
\usepackage{cite}
\usepackage{amsmath,amssymb,amsfonts}
\usepackage{algorithmic}
\usepackage{graphicx}
\usepackage{textcomp}
\usepackage{xcolor}
\usepackage{hyperref}

% Optional: For better figure handling
\usepackage{caption}
\usepackage{subcaption}

% Optional: For custom colors if needed
\definecolor{IEEEblue}{RGB}{0, 51, 153}
\hypersetup{
    colorlinks=true,
    linkcolor=IEEEblue,
    citecolor=IEEEblue,
    urlcolor=IEEEblue
}

% Title and Author Information
\title{Title of the Conference Paper}

\author{\IEEEauthorblockN{First Author}
\IEEEauthorblockA{\textit{Dept. of Computer Science} \\
\textit{University Name}\n
City, Country \\
email@example.com}
\and
\IEEEauthorblockN{Second Author}
\IEEEauthorblockA{\textit{Dept. of Electrical Engineering} \\
\textit{University Name}\n
City, Country \\
email@example.com}
}

\begin{document}

\section{Abstract}
% Briefly summarize the problem, methodology, key results, and conclusions. Keep it concise (usually 150-250 words).

\section{Introduction}
% Introduce the background, motivation, and contributions of the paper. Clearly state the problem being addressed and outline the structure of the paper.

\section{Related Work}
% Review existing literature and highlight how your work differs from or improves upon previous studies.

\section{Methodology}
% Describe the proposed approach, algorithms, models, or experimental setup in detail. Use equations and diagrams where necessary.

\section{Experiments and Results}
% Present the experimental setup, datasets, metrics, and results. Include tables and figures to support your findings. Analyze the results critically.

\section{Discussion}
% Interpret the results, discuss limitations, and explain the implications of your findings.

\section{Conclusion}
% Summarize the main contributions and suggest directions for future work.

\section{References}
% List all cited works in IEEE format. Use BibTeX if available.

\end{document}